\section{Exact finite inertia at the old bad places}
\label{sec:cubic-away-types}
This section keeps \(p=269\), \(K=\mathbb Q(\omega)\), the actual odd-weight
hypotheses and the ordinary GL(2) corner of AC-G78--AC-G79.  In particular,
all their original arithmetic and coefficient-ring hypotheses remain in
force.  It refines the old bad-place condition in that construction.
The new auxiliary place \(q\), the size of \(q\), and the per-curve
Hecke-support condition are separate.

\begin{proposition}[AC-G80-inertia: the actual killing condition supplies a split type]
\label{prop:G80-inertia}
Let \(v\nmid p\) be a place of a field \(F\supset K\), and let \(V\) be the
old two-dimensional endpoint representation obtained by the blockwise
untwist in AC-K29-point and AC-G68.  Write
\(\bar\gamma=\bar\Psi\,\bar\epsilon^{\,e}\).
Retain the actual old killing condition
\[
 g\in I_v,\quad \bar\Psi(g)=1
 \quad\Longrightarrow\quad V(g)=1.                 \tag{G80.1}
\]
Retain also the integral residual identification
\[
 \bar V|_{I_v}^{\rm ss}
       =\bar\Psi|_{I_v}\oplus\bar\Psi^{-1}|_{I_v}.  \tag{G80.2}
\]
Then \(V|_{I_v}\) is the direct sum of the unique prime-to-\(p\)
Teichmuller lifts of the two displayed characters.  In particular its
image is finite cyclic of order prime to \(p\), and its monodromy is zero.
The inertia characters extend to the local Weil group.

On the minimal ordinary GL(2) Hecke branch through the original classical
form \(f_0\), all arithmetic specializations used in AC-G78 have the same
finite inertia representation after restriction to \(F_v\), and have
zero monodromy there.  This conclusion permits a supercuspidal
representation over \(\mathbb Q_\ell\); the split assertion concerns its
restriction to \(F_v\).
\end{proposition}

\begin{proof}
Cyclotomic inertia is trivial away \(p\), and \(I_v\subset G_K\).  Thus
(G80.1) is exactly the restriction of the old killing condition: the
quadratic multiplier is one and the residual seed is the identity
precisely when \(\bar\Psi(g)=1\).
Put \(C=\operatorname{im}(\bar\Psi|_{I_v})\).  It is a finite subgroup of
the multiplicative group of a finite field of characteristic \(p\);
hence it is cyclic of order \(d\) prime to \(p\).  Condition (G80.1)
makes \(V|_{I_v}\) a representation of \(C\), rather than merely a
representation with a finite semisimplification.

After a finite unramified coefficient extension containing \(\mu_d\),
the idempotents
\[
 e_\xi=\frac1d\sum_{c\in C}\xi(c)^{-1}V(c)          \tag{G80.3}
\]
are integral orthogonal projectors on the stable lattice.  Their
reductions are the corresponding projectors for \(\bar V\).
Each summand is a direct summand of a free lattice, so its rank equals
its residual dimension.  Equation (G80.2) therefore gives exactly the
two claimed characters, with their stated multiplicities even if they
coincide.  There is no extra characteristic-zero character with zero
residual multiplicity.  Finite inertia implies \(N=0\).
The character \(\bar\Psi\) is the restriction of a global rational
character; conjugation by local Frobenius fixes it.  The same holds for
its unique Teichmuller lift.  Extend a character of inertia to the Weil
group by assigning a value to a Frobenius lift.  Each isotypic block is
therefore Frobenius-stable; Frobenius semisimplification only introduces
unramified twists inside a repeated-character block.

For the family statement use Saha, \emph{Purity for families of Galois
representations}, Theorem 6.2, printed pp.~901--902, on the minimal
ordinary Hecke branch through \(f_0\).  Its three alternatives describe
the generic local Weil--Deligne representation and every arithmetic
specialization.  The special representation alternative has nonzero
monodromy at every such point, so the preceding \(N=0\) computation at
\(f_0\) excludes it.  In each remaining alternative the finite inertia
representation is constant.  More explicitly, its matrices factor
through one finite inertia quotient.  Over an algebraic extension of the
coefficient fraction field its character values are sums of roots of
unity and are algebraic constants.  Choose the constant embedding
compatible with \(f_0\).  All characteristic-zero specializations preserve
these constants, and finite-group semisimplicity identifies their inertia
representations.  If necessary this is done after a finite coefficient
extension and on a component through \(f_0\); arithmetic points remain
dense there.  Restrict to \(I_v\) and apply the computation above.
At each actual old rational place the killing argument also applies
over \(K\), and the corresponding inertia subgroup has index at most
two in rational inertia.  Thus the original rational finite inertia
image has order prime to \(p\) as well; splitting these finite character
constants alone requires no ramified coefficient extension.
The analytic half-weight normalization, the cyclotomic twist, and the
fixed conductor-\(p\) quadratic character are unramified at \(v\), and
therefore do not alter this assertion.

Saha's theorem assumes generic absolute irreducibility of the classical
ordinary branch, supplied in its setup.  It does not assume residual
irreducibility.  Its purity and local-global compatibility inputs are
retained published inputs; this argument does not replace them.
\end{proof}

\begin{proposition}[AC-G80-type: an integral compact character selecting zero monodromy]
\label{prop:G80-type}
Let \(L\) be a nonarchimedean local field of residue characteristic
different from \(p\), and let
\(\phi_1,\phi_2,\phi_3:\mathcal O_L^\times\to\overline{\mathbb Q}_p^\times\)
be finite characters of order prime to \(p\).
There are an explicit compact open subgroup \(C_\phi\subset
\mathrm{GL}_3(L)\) and a rank-one character
\(\mu_\phi:C_\phi\to\mathcal O^\times\), of order prime to \(p\), with the
following property.  An irreducible smooth characteristic-zero
representation \(\pi\) contains \(\mu_\phi\) if and only if its
Frobenius-semisimple local Langlands parameter has
\[
 N=0,\qquad r|_I\simeq
 \operatorname{rec}(\phi_1)|_I\oplus
 \operatorname{rec}(\phi_2)|_I\oplus
 \operatorname{rec}(\phi_3)|_I.                    \tag{G80.4}
\]
Here the reciprocity convention is fixed on both sides; unramified
normalizing twists do not change (G80.4).

For any integral module at level \(\ker\mu_\phi\), the desired character
part is the image of the integral idempotent
\[
 e_{\mu_\phi}=\frac1{|D_\phi|}
       \sum_{a\in D_\phi}\mu_\phi(a)^{-1}R_a,\qquad
 D_\phi=C_\phi/\ker\mu_\phi.                       \tag{G80.5}
\]
This assertion makes no prime-to-\(p\) assumption on the index of an
Iwahori subgroup in a repeated-character block.
\end{proposition}

\begin{proof}
Permute the three coordinates to group equal characters together, and
write \(M=\prod_b\mathrm{GL}_{m_b}\) for the resulting Levi subgroup.
Set
\[
 \mu_M((g_b)_b)=\prod_b\phi_b(\det g_b)
 \quad(g_b\in\mathrm{GL}_{m_b}(\mathcal O_L)).
                                                               \tag{G80.6}
\]
Extend each \(\phi_b\) to \(L^\times\), for example by giving a uniformizer
value one.  The character of the diagonal unit torus is strongly
parabolic: its Weyl stabilizer is exactly the product of the symmetric
groups on these equal-character blocks, and (G80.6) extends it to \(M\).

Here is the precise compact subgroup.  Choose a Borel compatible with
\(M\).  For a root \(\alpha\), let \(c_\alpha\) be the smallest positive
integer for which the torus character is trivial on
\(\alpha^\vee(1+\mathfrak p_L^{c_\alpha})\).
Give a positive root depth \(\lfloor c_\alpha/2\rfloor\) and a negative
root depth \(\lceil c_\alpha/2\rceil\).  Let \(J\) be the group generated
by the diagonal units and these root subgroups, and put
\[
             C_\phi=J\,M(\mathcal O_L).           \tag{G80.7}
\]
Kamgarpour--Schedler, \emph{Ramified Satake Isomorphisms},
Lemmas 43--44 and Proposition 45, show that this is a compact open
subgroup and give its unique character \(\mu_\phi\), equal to \(\mu_M\)
on \(M(\mathcal O_L)\) and trivial on the complementary root groups.
The decomposition into these groups shows
\(\operatorname{im}\mu_\phi=\operatorname{im}\mu_M\).
Its order consequently divides the least common multiple of the orders
of the \(\phi_i\), so it is prime to \(p\).  All its values can be placed
in a finite unramified coefficient extension.  Convention 6 of that
paper excludes no residue characteristic for type \(A\).

Their Theorem 10 identifies
\[
 \mathrm{cInd}_{C_\phi}^{\mathrm{GL}_3(L)}\mu_\phi
 \simeq i_P^{\mathrm{GL}_3(L)}
       \left(\mathrm{cInd}_{M(\mathcal O_L)}^{M(L)}\mu_M\right).
                                                               \tag{G80.8}
\]
An unramified modulus convention only twists the right side
unramifiedly.  The Weyl stabilizer computation also verifies the exact
hypothesis of Roche, \emph{Parabolic Induction and the Bernstein
Decomposition}, Theorem 3.1: induction from this Levi's principal-series
Bernstein block is an equivalence onto the corresponding block for
\(\mathrm{GL}_3\).  This is precisely the application in
Kamgarpour--Schedler, Proposition 47.

Consequently an irreducible representation containing \(\mu_\phi\)
has the form \(i_P\sigma\), where each factor of \(\sigma\), after its
\(\phi_b\circ\det\) twist is removed, has a
\(\mathrm{GL}_{m_b}(\mathcal O_L)\)-fixed vector.
One can see the last assertion by applying the induction equivalence
and Frobenius reciprocity to (G80.8).
The unramified local Langlands correspondence identifies these spherical
factors with parameters of trivial inertia and zero monodromy.
Compatibility with parabolic induction now gives (G80.4).
Conversely, a parameter satisfying (G80.4) splits into its distinct
inertia-character blocks.  After the fixed character twist each is
unramified with zero monodromy, hence corresponds to a spherical
representation of its \(\mathrm{GL}_{m_b}\).  Distinct blocks have
different inertial characters, so the same induction equivalence gives
the irreducible representation \(i_P\sigma\); (G80.8) gives its required
compact character vector.

This uses the usual characteristic-zero local Langlands correspondence
and its unramified and parabolic compatibilities as published inputs.
It makes no assertion that a rational type projector is integral.
Integrality instead follows directly from
\(p\nmid|D_\phi|\): (G80.5) is an idempotent over \(\mathcal O\).
In particular, when two \(\phi_i\) coincide, (G80.7) contains the
\emph{full} maximal compact of that \(\mathrm{GL}_2\) block.
A character twist of Steinberg has no such fixed vector, although it has
an Iwahori-fixed vector.  The possibly \(p\)-divisible index
\([\mathrm{GL}_2(\mathcal O_L):I_2]\) is never inverted.
\end{proof}

\begin{theorem}[AC-G80: retention of every old finite inertial type]
\label{thm:G80}
Retain all actual hypotheses of AC-G78 and AC-G79, including their
specified ordinary corner, scalar-witness condition when needed,
and the host \(F\) of AC-G77.  At the old places away \(p\), retain the
old-point killing condition (G80.1), as proved by the original
blockwise old-point construction.
Then the actual automorphic representation supplied by AC-G79 can be
chosen with all its conclusions and, in addition, the following
simultaneous conditions:
\[
 \begin{split}
  N_v&=0,\\
  r|_{I_v}&\simeq
       \widehat{\bar\Psi}|_{I_v}\oplus 1\oplus
       \widehat{\bar\Psi}^{-1}|_{I_v}
 \end{split}
 \qquad\text{at every old }v\nmid p.               \tag{G80.9}
\]
Here \(r=\epsilon r_\sigma\) is the unitary normalization of AC-G79;
the hats denote the unique prime-to-\(p\) lifts determined by the old
point.  At the remaining unchanged unramified places the representation
is unramified.  At \(p\) it has exactly the ordered flag of AC-G79,
and at the new \(q\) it has the finite dihedral type and zero monodromy
of AC-G78-type.  Its global rank-three Galois representation remains
absolutely irreducible.

This is an existence assertion in the stated odd-weight construction.
It does not assert the per-curve containment of AC-G61, a quantitative
bound for \(q\), or a new complete \(\mathrm{IsSeedAlt}\) theorem.
\end{theorem}

\begin{proof}
\textbf{The dense family has the desired old types.}
Include all old ramified primes in the finite set used to construct
\(F\) in AC-G77.  They split relative to \(F/F^+\), so their unitary
groups are the local general linear groups.
Proposition~\ref{prop:G80-inertia} and Saha's theorem show that the
dense classical GL(2) specializations used in AC-G78 have the fixed
finite split inertia after base change to \(F\).
The fixed cubic character \(\chi_{269}\) is unramified at these places.
At primes away \(3\cdot269\) this follows from the Kummer equation.
At \(3\), use \(269\equiv-1\pmod9\): \(-269\in1+9\mathbb Z_3\),
and \(\exp(\log(-269)/3)\in1+3\mathbb Z_3\) has cube \(-269\).
Its negative is a cube root of \(269\) in \(\mathbb Q_3\), so the
character is trivial on \(G_{K_3}\).
The fixed quadratic conductor-\(p\) character is also unramified here.  Adding the
middle line gives exactly the three characters in (G80.9).
Use the reciprocity convention of the retained local-global
compatibility to define \(\phi_v\) on local units.
Proposition~\ref{prop:G80-type} gives a compact pair
\((C_v,\mu_v)\) containing each of these classical local representations.

\textbf{Integral ordinary modules with these types.}
Take the actual tame level \(\prod_v\ker\mu_v\) at the old bad places,
hyperspecial level at the unchanged good places, and the same
sufficiently small auxiliary level used in AC-G72 and AC-G78.
Let \(D=\prod_v C_v/\ker\mu_v\).
This finite group has order prime to \(p\).
From the start retain the auxiliary-neat-level projector
\(e_{\rm aux}\) of AC-G72: choose its auxiliary prime outside the old
bad set and \(p\), with the finite maximal-compact quotient of order
prime to \(p\).  The dense family is unramified there.
Its integral averaging projector selects the maximal-compact invariant
part already on each fine-level module.  On these modules apply
\(e_{\rm old}=\prod_v e_{\mu_v}\), and throughout use the combined
projector \(e_{\rm aux}e_{\rm old}e_{\delta_e}\).
These commuting projectors commute with the coefficient-prime torus, its ordinary operator,
depth inclusions and traces, all Hecke operators outside the retained
bad set, and the full finite group at the new \(q\).
These are actions at different places.
The old-place and auxiliary projectors also commute with the finite
tame torus projector of AC-G79.

The sufficiently small double-coset construction gives the finite-depth
ordinary module as a free module over the local group ring of the
coefficient-prime depth group.  Taking these commuting integral
idempotents gives a projective, hence free, summand over that local
ring.  The trace base-change isomorphism and the identity
\[
 I_{d,d+1}\operatorname{Tr}_{d+1,d}(m)
       =\sum_{a\in\ker(D_{d+1}\to D_d)}a\,m        \tag{G80.10}
\]
restrict to this summand.
Compatible bases and the twisted norm argument of AC-G78-module
therefore give a finite free Iwasawa module with \emph{actual}
specializations
\[
 e_{\rm aux}e_{\rm old}e_{\delta_e}H_\lambda                 \tag{G80.11}
\]
for every continuous weight used in that argument.
No projection from an Iwahori level by a Weyl-group average occurs.

\textbf{The full-chain fibre and the four ordinary tests.}
Every dense classical endoscopic specialization has a vector in
(G80.11).  Repeat the finite-Hecke-algebra construction of AC-G78 on
this finite free module.  Every actual operator relation vanishes at
these dense classical points, so its desired pseudotrace value
vanishes in the branch ring.  This defines the same branch-valued
Hecke character.  The analytic root argument chooses the same kind of
auxiliary \(q\).  The Cayley--Hamilton polynomial-law argument,
valid after arbitrary base change, makes the specialized character a
character of the \emph{actual} Hecke image on (G80.11).

The maps of the elliptic \(q\)-complex in AC-G75 act only at \(q\) and
thus commute with both the old-place and auxiliary projectors.
Apply its generic cohomological concentration at the actual fine level,
with \(\ker\mu_v\) at the old places, and then take this combined
integral direct summand.
The involution \(\Theta\) sends
\[
 (C_v,\mu_v)\ \longmapsto\
 (\Theta(C_v),\,\mu_v\circ\Theta^{-1});             \tag{G80.12}
\]
the pullback of functions intertwines their exact character parts.
The transformed character again has prime-to-\(p\) image.
At the auxiliary neat place, transpose-inverse preserves both the
maximal compact and its principal congruence subgroup, so it preserves
the trivial-character projector \(e_{\rm aux}\).
Thus the fourth ordinary test is performed at these transformed
compact pairs, just as it uses the transformed coefficient-prime
character in AC-G79.  The other three tests use the original pairs.
The no-Steinberg premise is transported by this isomorphism, so the
same four-test contradiction gives an actual ordinary Steinberg
constituent at \(q\) in the original fixed weight and the original
old-type summand.

\textbf{The finite \(q\)-type and all local conclusions.}
At full \(q\)-level the fixed-weight lattice is still projective over
\(\mathcal O[\mathrm{GL}_3(k_q)]\): the auxiliary, old-type and
coefficient-prime projectors commute with that action.
The residual genericity witness of AC-G78 is unchanged.
Every constituent is therefore generic at split places.  Its spherical
\(q\)-constituents contribute zero to
\(m_{\mathrm{St}_3}-m_1\), its nonspherical constituents contribute
nonnegatively, and the newly supplied Steinberg contributes positively.
The integral-projective Brauer calculation of AC-G72 and AC-G78-type
supplies the actual finite dihedral type inside the same auxiliary, old-type and
coefficient-prime summand.
The local and global arguments of AC-G72-local and AC-G78-type give
zero monodromy at \(q\) and absolute global irreducibility.

At an old place the compact character vector gives (G80.4), hence
(G80.9), by Proposition~\ref{prop:G80-type}.
Here \(v\nmid p\), so full local-global compatibility identifies this
with the Galois Weil--Deligne parameter.  When \(N=0\), the original
inertia action itself is the finite inertia action of that parameter;
this is an assertion about actual inertia, not only its
semisimplification.  The coefficient-prime flag survives because its
projector was retained throughout.  Unchanged hyperspecial factors
give unramified parameters.  At the auxiliary neat place the representation is already in the
image of \(e_{\rm aux}\), which was retained in the universal module,
all four tests and the projective \(q\)-lattice.  It therefore has a
maximal-compact fixed vector there and is unramified.  No averaging
after the extraction of \(\tau_\theta\) is used.
\end{proof}

\paragraph{Exact source and dictionary scope.}
The new published inputs are Saha, Theorem 6.2 and its complete printed
proof (AIF 67 (2017), pp.~901--902); Kamgarpour--Schedler,
Convention 6, Theorem 10, Lemmas 43--44, Proposition 45 and
Proposition 47, with their applied printed proofs (Documenta Math.~18
(2013), pp.~1275--1300); and Roche, Theorem 3.1 with its proof in
Sections 3.1--3.5 (Compositio Math.~134 (2002), pp.~123--126).
The deeper type, purity and local Langlands theorems cited by these
papers remain imported results.
The old killing condition is the one already used and checked in
AC-K29-point and AC-G68.  Its precise pinned Lean relative path and a
fresh byte-level replay are deferred during the executor outage;
no new path is asserted from memory.
Geometrically, the finite character quotient is a finite etale
coefficient cover and (G80.5) is its prime-to-\(p\) character summand,
in the already audited SGA~1, Expose~V, Sections~4 and~7 framework.
No new archival Grothendieck cote or geometric realization theorem is
claimed.  The new \(q\)-type is finite but need not be diagonal over the
present host; applying a further cover, including AC-G73 when its
hypotheses are met, is a separate operation.
